%
% terminal_lifecycle_cache_leak_buggy.tex
%
% A Z specification of the lux terminal session lifecycle across the
% three processes that touch it: luxd-pty (owns shells), the terminal
% helper (the luxd-pty client that owns the Hub scenes and frames) and
% the Display (holds one replica per scene element).
%
% Design record: docs/architecture/terminal-lifecycle.md
% Source of record (today's code the buggy variants mirror):
%   src/punt_lux/ptyd/core.py              SessionCore.acquire
%   src/punt_lux/display/terminal/remote_pool.py   RemotePool.get
%   src/punt_lux/display/terminal/attachment_requests.py
%
% Type-check:  fuzz -t docs/terminal_lifecycle_cache_leak_buggy.tex
% Model-check (see Section "Verification"):
%   probcli docs/terminal_lifecycle_cache_leak_buggy.tex -model_check -p DEFAULT_SETSIZE 2
%
\documentclass[a4paper,10pt,fleqn]{article}
\usepackage[margin=1in]{geometry}
\usepackage{fuzz}

\begin{document}

\title{The lux Terminal Session Lifecycle: a Z Specification}
\author{Formal model of exit, resume and restart}
\date{November 2026}
\maketitle

\section{Introduction}

A terminal session is born in the daemon, shown by the helper as a Hub
scene (whose frame dies with it), and painted by a Display replica.
Three defects recurred in one week because no one document said who
removes what, and when.  The user-visible form of all three: a frame
that is on screen but can never attach again.

This model makes removal an obligation.  Every visible frame is either
attachable, or the helper has been told and owes its removal.

\section{Basic Types}

\begin{zed}
  [NAME, SHELL]
\end{zed}

\texttt{NAME} is the session \emph{id} on one host (the design's daemon-assigned, never-reused id; labels are not modelled).  \texttt{SHELL} is the
identity of one spawned program: the daemon's generation number.  A name
may be reused by a later shell; a shell identity is never reused while
anything still refers to it.

\begin{zed}
  DSTATE ::= running | exited \\
  RPHASE ::= ready | gone \\
  SYNC ::= synced | unsynced
\end{zed}

\texttt{exited} is the daemon's draining state: the program has ended
and the viewers are being told.  \texttt{gone} is a replica that ended
(exited or failed to attach).  \texttt{unsynced} means the helper may
hold scenes or have missed events that no longer match the daemon.

\section{State}

\begin{schema}{Sys}
  daemon : NAME \pfun SHELL \\
  dstate : NAME \pfun DSTATE \\
  scene : NAME \pfun SHELL \\
  replica : (NAME \cross SHELL) \pfun RPHASE \\
  bound : (NAME \cross SHELL) \pfun SHELL \\
  createdEv, endedEv : NAME \rel SHELL \\
  used : \power NAME \\
  sync : SYNC
\where
  \dom dstate = \dom daemon \\
  daemon \in NAME \pinj SHELL \\
  \dom bound = \dom replica
\end{schema}

\texttt{daemon} is the single authoritative map from a name to its live
shell on this host; being a function it gives at most one shell per
name, and being injective at most one name per shell.  \texttt{scene}
is what the Hub shows: the element's name and the shell it expects.
\texttt{replica} is the Display's cache, keyed by name \emph{and} shell.
\texttt{bound} records which shell a replica's stream actually comes
from.  \texttt{createdEv} and \texttt{endedEv} are daemon notices in
flight to the helper.

\begin{schema}{Init}
  Sys'
\where
  daemon' = \emptyset \\
  dstate' = \emptyset \\
  scene' = \emptyset \\
  replica' = \emptyset \\
  bound' = \emptyset \\
  createdEv' = \emptyset \\
  endedEv' = \emptyset \\
  used' = \emptyset \\
  sync' = synced
\end{schema}

\section{Daemon Operations}

\begin{schema}{DaemonNew}
  \Delta Sys \\
  n? : NAME \\
  s? : SHELL
\where
  n? \notin \dom daemon \\
  s? \notin \ran daemon \cup \ran scene \cup \ran createdEv \cup
    \ran endedEv \cup \ran (\dom replica) \\
  daemon' = daemon \cup \{ n? \mapsto s? \} \\
  dstate' = dstate \cup \{ n? \mapsto running \} \\
  createdEv' = createdEv \cup \{ (n? , s?) \} \\
  used' = used \cup \{ n? \} \\
  scene' = scene \\
  replica' = replica \\
  bound' = bound \\
  endedEv' = endedEv \\
  sync' = sync
\end{schema}

Creation always announces itself, so the helper shows a scene whoever
asked: the menu, the CLI, another client.

\begin{schema}{ShellExit}
  \Delta Sys \\
  n? : NAME
\where
  used' = used \\
  n? \in \dom daemon \\
  dstate~n? = running \\
  dstate' = dstate \oplus \{ n? \mapsto exited \} \\
  endedEv' = endedEv \cup \{ (n? , daemon~n?) \} \\
  daemon' = daemon \\
  scene' = scene \\
  replica' = replica \\
  bound' = bound \\
  createdEv' = createdEv \\
  sync' = sync
\end{schema}

Exit and an explicit kill are the same transition: the program ends and
the daemon announces it.

\begin{schema}{DaemonRemove}
  \Delta Sys \\
  n? : NAME
\where
  used' = used \\
  n? \in \dom daemon \\
  dstate~n? = exited \\
  (n? , daemon~n?) \in \dom replica \implies
    replica (n? , daemon~n?) \neq ready \\
  daemon' = \{ n? \} \ndres daemon \\
  dstate' = \{ n? \} \ndres dstate \\
  scene' = scene \\
  replica' = replica \\
  bound' = bound \\
  createdEv' = createdEv \\
  endedEv' = endedEv \\
  sync' = sync
\end{schema}

The session is removed only once no viewer is still being served from
it: every attached replica has been told.

\begin{schema}{DaemonRestart}
  \Delta Sys
\where
  used' = used \\
  daemon' = \emptyset \\
  dstate' = \emptyset \\
  createdEv' = \emptyset \\
  endedEv' = \emptyset \\
  sync' = unsynced \\
  scene' = scene \\
  replica' = replica \\
  bound' = bound
\end{schema}

A daemon restart ends every program; persistence across it is a
non-goal.  The helper is the daemon's child and restarts with it.

\section{Helper Operations}

\begin{schema}{HelperCreated}
  \Delta Sys \\
  n? : NAME \\
  s? : SHELL
\where
  used' = used \\
  (n? , s?) \in createdEv \\
  n? \in \dom daemon \\
  daemon~n? = s? \\
  dstate~n? = running \\
  scene' = scene \oplus \{ n? \mapsto s? \} \\
  createdEv' = createdEv \setminus \{ (n? , s?) \} \\
  daemon' = daemon \\
  dstate' = dstate \\
  replica' = replica \\
  bound' = bound \\
  endedEv' = endedEv \\
  sync' = sync
\end{schema}

\begin{schema}{HelperCreatedStale}
  \Delta Sys \\
  n? : NAME \\
  s? : SHELL
\where
  used' = used \\
  (n? , s?) \in createdEv \\
  \lnot (n? \in \dom daemon \land daemon~n? = s?
    \land dstate~n? = running) \\
  createdEv' = createdEv \setminus \{ (n? , s?) \} \\
  daemon' = daemon \\
  dstate' = dstate \\
  scene' = scene \\
  replica' = replica \\
  bound' = bound \\
  endedEv' = endedEv \\
  sync' = sync
\end{schema}

A creation notice that arrives after its shell has already ended is
dropped; the ended notice handles the rest.

\begin{schema}{HelperEnded}
  \Delta Sys \\
  n? : NAME \\
  s? : SHELL
\where
  used' = used \\
  (n? , s?) \in endedEv \\
  endedEv' = endedEv \setminus \{ (n? , s?) \} \\
  (n? \in \dom scene \land scene~n? = s?) \implies
    scene' = \{ n? \} \ndres scene \\
  \lnot (n? \in \dom scene \land scene~n? = s?) \implies
    scene' = scene \\
  daemon' = daemon \\
  dstate' = dstate \\
  createdEv' = createdEv \\
  sync' = sync \\
  replica' = replica \\
  bound' = bound
\end{schema}

The helper clears the scene only if it still names the ended shell, so
a late notice never removes a newer terminal that reuses the name.

\begin{schema}{HelperLoseEvents}
  \Delta Sys
\where
  used' = used \\
  createdEv' = \emptyset \\
  endedEv' = \emptyset \\
  sync' = unsynced \\
  daemon' = daemon \\
  dstate' = dstate \\
  scene' = scene \\
  replica' = replica \\
  bound' = bound
\end{schema}

The helper crashes or its watch connection drops.  Notices in flight are
lost; the next step must reconcile.

\begin{schema}{HelperSync}
  \Delta Sys
\where
  used' = used \\
  sync = unsynced \\
  scene' = \{ n : NAME | n \in \dom daemon \land dstate~n = running
    @ n \mapsto daemon~n \} \\
  createdEv' = \emptyset \\
  endedEv' = \emptyset \\
  sync' = synced \\
  daemon' = daemon \\
  dstate' = dstate \\
  replica' = replica \\
  bound' = bound
\end{schema}

Reconciliation replaces the helper's scenes with exactly the running
sessions the daemon lists.

\section{Display Operations}

\begin{schema}{ReplicaAttach}
  \Delta Sys \\
  n? : NAME
\where
  used' = used \\
  n? \in \dom scene \\
  (n? , scene~n?) \notin \dom replica \\
  n? \in \dom daemon \\
  daemon~n? = scene~n? \\
  replica' = replica \cup \{ (n? , scene~n?) \mapsto ready \} \\
  bound' = bound \cup \{ (n? , scene~n?) \mapsto daemon~n? \} \\
  daemon' = daemon \\
  dstate' = dstate \\
  scene' = scene \\
  createdEv' = createdEv \\
  endedEv' = endedEv \\
  sync' = sync
\end{schema}

Resume names the shell the scene expects and attaches only to that one.

\begin{schema}{ReplicaRefused}
  \Delta Sys \\
  n? : NAME
\where
  used' = used \\
  n? \in \dom scene \\
  (n? , scene~n?) \notin \dom replica \\
  \lnot (n? \in \dom daemon \land daemon~n? = scene~n?) \\
  replica' = replica \cup \{ (n? , scene~n?) \mapsto gone \} \\
  bound' = bound \cup \{ (n? , scene~n?) \mapsto scene~n? \} \\
  daemon' = daemon \\
  dstate' = dstate \\
  scene' = scene \\
  createdEv' = createdEv \\
  endedEv' = endedEv \\
  sync' = sync
\end{schema}

\begin{schema}{ReplicaEnds}
  \Delta Sys \\
  n? : NAME \\
  s? : SHELL
\where
  used' = used \\
  (n? , s?) \in \dom replica \\
  replica (n? , s?) = ready \\
  \lnot (n? \in \dom daemon \land daemon~n? = s? \land dstate~n? = running) \\
  replica' = replica \oplus \{ (n? , s?) \mapsto gone \} \\
  bound' = bound \\
  daemon' = daemon \\
  dstate' = dstate \\
  scene' = scene \\
  createdEv' = createdEv \\
  endedEv' = endedEv \\
  sync' = sync
\end{schema}

The replica learns of the end from the stream's close notice or from the
lost connection.



A replica no scene element names is dropped.  Today's pool never does
this.

\section{Invariants}
\label{sec:inv}

None is placed in \texttt{Sys}: a predicate there would disable the very
transition that breaks it.  Each is checked as an unreachable goal.

\paragraph{I1 --- visible implies attachable or being removed.}
Every scene names a running shell on the daemon, or a removal notice for
it is in flight, or the helper is known to be out of sync:
\[
  \forall n : \dom scene @ sync = unsynced \lor
    (n \in \dom daemon \land daemon~n = scene~n \land dstate~n = running)
    \lor (n , scene~n) \in endedEv
\]

\paragraph{I2 --- resume never swaps.}
Every replica streams from the shell it names:
\[
  \forall k : \dom bound @ bound~k = \mathrm{second}~k
\]

\paragraph{I3 --- one live shell per name.}
Holds by construction (\texttt{daemon} is an injective function), and is
checked as: no two names share a shell.

\paragraph{I4 --- replicas learn of a lost session (transient, not a safety property).}
ProB found a state with a \texttt{ready} replica of a name the daemon
no longer holds: it arises only after \texttt{DaemonRestart}, before
\texttt{ReplicaEnds} fires.  It is a legitimate transient, so it is
\emph{not} an invariant.  What must hold is that \texttt{DaemonRemove}
(the exit path) never leaves one behind, which its guard enforces, and
that \texttt{ReplicaEnds} stays enabled in that state (no deadlock).

\paragraph{I6 --- no stale replica shadows a live scene.}
The Display caches one replica per session.  A scene that names a running
shell must not coexist with a replica of the same name for a different
shell, because that cached replica would be served instead (today's
immortal \texttt{RemotePool} cache, keyed by name):
\[
  \forall n : \dom scene @ \lnot (n \in \dom daemon \land
    daemon~n = scene~n \land dstate~n = running \land
    \exists s : SHELL @ (n , s) \in \dom replica \land s \neq scene~n)
\]
Ids are never reused, so in the intact design this is unreachable; the
\texttt{used} set models that guarantee.

\paragraph{I5 --- deadlock freedom.}
Some operation is always enabled.

\section{Verification}
\label{sec:verify}

\begin{verbatim}
  fuzz -t docs/terminal_lifecycle_cache_leak_buggy.tex
  probcli docs/terminal_lifecycle_cache_leak_buggy.tex -model_check -p DEFAULT_SETSIZE 2
  # I1 .. I4: each goal must report "not found"
\end{verbatim}

\section{Fidelity}

Two buggy variants of this file exist beside it:
\texttt{terminal\_lifecycle\_exit\_leak\_buggy.tex} drops the ended
notice and the removal (today: exit leaves a frame that can never attach
again), and \texttt{terminal\_lifecycle\_swap\_buggy.tex} lets a resume
attach by name alone (the exit-respawn swap).  Each must make a goal
reachable that the intact model cannot reach.

\end{document}
